\section{Cryptographic Audit Trail}
\label{sec:crypto}

A redacted document is only trustworthy if a verifier
can confirm \emph{independently} that redaction
occurred, when it occurred, and that the process was
executed correctly.
Manual audit logs are mutable and unverifiable;
centralised databases require trust in their operator.
Omni-Shield addresses this through two complementary
cryptographic mechanisms: a blockchain-anchored audit
ledger that provides tamper-evident provenance, and
a Groth16 zero-knowledge proof that attests to the
correctness of the redaction without disclosing
what was redacted.
Both operate asynchronously --- completing after the
redacted document is returned to the user --- and
impose no latency penalty on the interactive
redaction workflow.

% ─────────────────────────────────────────────────────────
\subsection{Blockchain Audit Ledger}
% ─────────────────────────────────────────────────────────

\paragraph{Smart contract design.}
The \texttt{AuditLog.sol} Solidity contract maintains
two on-chain data structures per record:
a \texttt{mapping(bytes32\,$\Rightarrow$\,Record)}
for O(1) hash lookup, and a per-owner
\texttt{mapping(address\,$\Rightarrow$\,bytes32[])}
index that allows a wallet holder to enumerate all
documents they have processed.
Each \texttt{Record} stores the anchor hash,
redaction count, block timestamp, and document type.
The \texttt{addRecord()} function enforces an
idempotency check --- anchoring the same hash twice
is silently skipped --- and emits a
\texttt{RecordAdded} event for off-chain indexing.
The \texttt{verify()} view function is gas-free
(read-only), enabling third-party verification
without transaction cost.

\paragraph{Salted hash commitment.}
A naive implementation would anchor
\texttt{SHA256(document\_bytes)} directly.
This introduces a linkability vulnerability: identical
documents processed multiple times produce identical
on-chain hashes, allowing an observer monitoring the
blockchain to detect repeated processing of the same
document.
Analysis of our audit log reveals 260 duplicate
hashes among 424 records, all arising from test
documents processed repeatedly during evaluation.

We address this by computing a salted anchor hash:
\begin{equation}
  h_{\text{anchor}} = \text{SHA256}(h_{\text{doc}}
    \,\|\, \text{owner} \,\|\, \tau)
  \label{eq:salt}
\end{equation}
where $h_{\text{doc}} = \text{SHA256(document\_bytes)}$
is the original document hash, $\text{owner}$ is the
wallet address, and $\tau$ is the Unix timestamp at
anchor time.
The original $h_{\text{doc}}$ is preserved in the
audit JSON for user-side verification (the user can
recompute it from their document);
$h_{\text{anchor}}$ is what appears on-chain.
The verification endpoint performs a two-step lookup:
given $h_{\text{doc}}$ from the user, it retrieves
$h_{\text{anchor}}$ from the stored audit record and
queries the blockchain for that value.

Equation~\ref{eq:salt} ensures that even if two users
process the same document, their on-chain anchors are
distinct.
An adversary who observes the blockchain learns only
that \emph{some} document was processed by
\emph{some} wallet at \emph{some} time --- no
document identity information is recoverable from
the on-chain record alone.

\paragraph{Deployment and gas cost.}
The system currently deploys to a local Ganache
Ethereum testnet~\cite{ganache} for development.
Each \texttt{addRecord()} call consumes 264,333 gas
units.
At Ethereum mainnet prices (\$2,500 per ETH, 20\,Gwei
gas price), this corresponds to approximately
\$13.22 per anchor --- prohibitive for high-volume
use.
Deployment on Polygon L2~\cite{polygon2021} reduces
this to approximately \$0.006 per anchor at 30\,Gwei
and \$0.80 per MATIC, making production deployment
economically viable.
The \texttt{AuditLog.sol} contract is EVM-compatible
and requires only an environment variable change to
migrate between networks.

\paragraph{Verification performance.}
The \texttt{verify()} view function resolves in
1.2\,ms mean (measured over five queries to the local
node), reflecting the O(1) mapping lookup.
This is fast enough to be embedded in document
workflows without perceptible latency --- a verifier
can confirm a redaction record in under 2\,ms,
faster than a typical HTTP round-trip to a remote
audit service.

% ─────────────────────────────────────────────────────────
\subsection{Zero-Knowledge Proof of Redaction}
% ─────────────────────────────────────────────────────────

\paragraph{Motivation.}
The blockchain record proves that \emph{someone}
anchored a hash at a specific time.
It does not prove that the redaction was performed
correctly --- a malicious actor could anchor a hash
corresponding to an unredacted document.
A zero-knowledge proof closes this gap: it attests
that the prover knows a set of redacted values
consistent with the public commitment, without
revealing the values themselves.

\paragraph{Circuit design.}
We implement the ZK circuit in ZoKrates~\cite{eberhardt2018zokrates}
using the Groth16 proving system over the
BN128 elliptic curve.
The circuit proves the following statement:

\begin{quote}
\textit{I know a private witness
$\mathbf{w} = \{h_{\text{items}}\}$ such that
$\text{Commit}(h_{\text{doc}},\, h_{\text{items}}) =
c_{\text{public}}$, where $c_{\text{public}}$ is the
publicly known commitment, $h_{\text{doc}}$ is the
public document hash, and $h_{\text{items}}$ is the
SHA-256 hash of the concatenated redacted values
(known only to the prover).}
\end{quote}

The commitment function uses Pedersen hashing
(available in the ZoKrates standard library) over
the concatenation of the document hash bits and
the items hash bits:
\begin{equation}
  c = \text{PedersenHash}(
        h_{\text{doc}} \,\|\, h_{\text{items}})
  \label{eq:commitment}
\end{equation}
The circuit takes as \emph{public inputs}:
the document hash $h_{\text{doc}}$ (as eight
32-bit unsigned integers), the redaction count,
and the commitment $c$ (two 32-bit words).
The circuit takes as \emph{private input}
(the witness): $h_{\text{items}}$ (eight 32-bit
unsigned integers).
The private input never leaves the proving system
and is not included in the proof.

\paragraph{Proof generation and verification.}
Proof generation via ZoKrates \texttt{generate-proof}
takes 20.3\,s mean on the target hardware.
This is executed asynchronously in a background
thread immediately after the redacted document is
saved; the user receives their output before proof
generation completes.
The proof becomes available in the audit record
within approximately 25\,s of the redaction request.

Proof verification uses ZoKrates \texttt{verify}
against the on-disk verification key generated
during the one-time trusted setup.
Verification takes 22.1\,ms mean ---
fast enough for interactive verification workflows.

Groth16 proofs are constant size (128 bytes for the
proof object, three field elements for public inputs)
regardless of the number of redacted items or the
document size.
This makes proof storage and transmission trivially
cheap.

\paragraph{Trusted setup.}
Groth16 requires a one-time trusted setup ceremony
to generate the proving key (used by the prover) and
verification key (used by verifiers).
The current implementation uses a locally generated
setup, which introduces a trust assumption: the party
who ran the setup could theoretically generate false
proofs for arbitrary statements.
A production deployment should replace the local
setup with a publicly verifiable multi-party
computation (MPC) ceremony, or migrate to a
transparent proof system (PLONK, STARKs) that
eliminates the trusted setup requirement entirely.
We identify this as a priority for future work
(Section~\ref{sec:limitations}).

\paragraph{What the proof guarantees.}
The ZK proof provides a \emph{selective disclosure}
property: an auditor can verify that a specific
document was redacted and that the redaction
committed to a specific (undisclosed) set of PII
values, without learning what those values were.
This is strictly stronger than the blockchain
record alone, which only proves that \emph{a}
hash was anchored.
Combined, the two mechanisms provide:

\begin{itemize}
  \item \textbf{Tamper evidence}: the document
        hash is immutable on-chain; any alteration
        to the redacted output invalidates the record.
  \item \textbf{Provenance}: the owner wallet
        address and timestamp are permanently
        associated with the record.
  \item \textbf{Redaction correctness}: the ZK
        proof attests that the committed PII values
        are consistent with the claimed document.
  \item \textbf{Privacy}: neither the blockchain
        record nor the proof reveals which values
        were redacted or their content.
\end{itemize}

% ─────────────────────────────────────────────────────────
\subsection{Audit Record Structure}
% ─────────────────────────────────────────────────────────

Every completed redaction writes a
\texttt{.audit.json} sidecar file alongside the
output document.
The sidecar contains:

\begin{itemize}
  \item \texttt{sha256\_hex}: original document hash
        (user-verifiable by re-hashing the input)
  \item \texttt{anchor\_hash}: salted on-chain hash
  \item \texttt{anchor\_salt}: the salt string used
        (required to recompute \texttt{anchor\_hash}
        for verification)
  \item \texttt{redaction\_count}: number of PII
        items redacted
  \item \texttt{coverage\_pct}: percentage of document
        area or text volume redacted
  \item \texttt{pii\_filter}: the category selection
        active at redaction time
  \item \texttt{redactions}: list of redacted items
        with masked values and bounding boxes
  \item \texttt{zk\_proof}: Groth16 proof object,
        public inputs, and commitment
  \item \texttt{timestamp}: Unix timestamp of
        redaction
  \item \texttt{doc\_type}: inferred document type
\end{itemize}

The sidecar is written synchronously before
blockchain anchoring and proof generation begin,
ensuring it is always available even if the
asynchronous background tasks fail.
A \texttt{GET /api/audit/\{sha256\}} endpoint serves
sidecar data by hash prefix, enabling the frontend
to retrieve and display the full audit record
for any previously processed document.

% ─────────────────────────────────────────────────────────
\subsection{Threat Model}
% ─────────────────────────────────────────────────────────

We consider three classes of adversary.

\textbf{Passive blockchain observer.}
An adversary who can read all on-chain transactions
learns: which wallet addresses have processed
documents, how many items were redacted, and when.
They learn nothing about document content, PII
values, or which documents were processed (because
anchor hashes are salted).
This is a substantial improvement over cloud-based
alternatives, where the service provider learns
the document content in full.

\textbf{Malicious prover.}
An adversary who controls the proving system
and wishes to generate a false proof (claiming a
document was redacted when it was not) cannot do
so without knowledge of the trusted setup trapdoor.
Under the assumptions of the Groth16 proof system
(computational soundness under the knowledge
of exponent assumption), such forgeries are
computationally infeasible.
The local trusted setup weakens this guarantee:
if the setup party is malicious, false proofs
are possible.
The MPC ceremony mitigation eliminates this
assumption.

\textbf{Compromised client.}
An adversary who compromises the user's device
after redaction can access the audit sidecar,
including the anchor salt and redaction list.
They cannot alter the on-chain record retroactively.
The sidecar should be treated with the same access
controls as the original document.

\paragraph{Out of scope.}
We do not address: network-level attacks on the
local Ganache node (protected by firewall in
production), adversarial inputs designed to evade
the VLM, or side-channel attacks on the GPU
inference process.
These are identified as areas for future security
analysis.

% ─────────────────────────────────────────────────────────
\subsection{Energy and Cost Analysis}
% ─────────────────────────────────────────────────────────

\begin{table*}[h]\centering
\caption{Energy Consumption and Carbon Footprint per 1,000 Redactions}
\label{tab:energy}
\begin{tabular}{llrrr}
\toprule
System & Type & kWh/1K & CO$_2$e (kg) & Grid assumed \\
\midrule
Omni-Shield (image) & Local edge  & 0.289 & 0.237 & India (0.82 kg/kWh) \\
Omni-Shield (image) & Local edge  & 0.289 & 0.113 & US avg (0.39 kg/kWh) \\
Omni-Shield (image) & Local edge  & 0.289 & 0.023 & Renewable (0.08 kg/kWh) \\
Omni-Shield (PDF)   & Local edge  & 0.001 & 0.001 & India \\
AWS Comprehend      & Cloud NLP   & 0.200 & 0.030 & AWS renewable \\
Google Cloud DLP    & Cloud NLP   & 0.200 & 0.016 & GCP renewable \\
GPT-4 Vision        & Cloud LLM   & 1.000 & 0.150 & AWS renewable \\
\bottomrule
\end{tabular}
\footnotesize{Omni-Shield power draw: 100W total (RTX 5060 75W + 25W overhead).
Cloud figures estimated from Patterson et al. 2022.}
\end{table*}


Table~\ref{tab:energy} reports energy consumption
and estimated carbon footprint for Omni-Shield and
comparable cloud alternatives.
At 100\,W total system draw (RTX~5060 75\,W plus
25\,W overhead), Omni-Shield consumes 0.289\,kWh
per 1,000 image redactions --- comparable to AWS
Comprehend (estimated 0.200\,kWh/1K) and
substantially lower than GPT-4~Vision
(estimated 1.000\,kWh/1K).

Carbon footprint is highly sensitive to grid carbon
intensity.
On the Indian grid (0.82\,kg CO$_2$/kWh),
Omni-Shield produces 0.237\,kg CO$_2$e per 1,000
redactions.
On a renewable electricity tariff (0.08\,kg/kWh),
this falls to 0.023\,kg --- 1.3$\times$ lower than
AWS Comprehend on its renewable energy commitment.

The additional energy cost of local-first processing
relative to a cloud NLP API is 0.089\,kWh per 1,000
documents, equivalent to approximately \rupee0.72 at
Indian electricity retail rates.
This is the quantified \emph{privacy premium}: the
energy overhead of the guarantee that sensitive
identity documents never leave the device.
Whether this premium is acceptable is a
context-dependent policy decision; we report it
as a measured quantity rather than an argument for
or against local deployment.
